\subsection{Which service policies admit a regional material solve?}
Removing the graph service cap lets newly ignited sites donate immediately and invalidates the old fixed-donor argument. Under this policy we take the union of connected components of the native donor-neighbor graph that intersect the queried shell. During a certified emission prefix no donor outside this union can deposit graph heat inside. Induct on local cold and burning states, keeping all new local donors in native order and performing full global gas updates with the same proven emitters. This gives the correct local trajectory under full donor service. If the domain spans the asset, we use ordinary replay.

The trained playground has four components. The two queried unions contain 30,459 and 9,810 of 80,023 sites. Both edits reached the shell at step $+2$ with 3,654 and 4,920 eligible sites. The component method passed the native state audit, including 1,819 and 987 new local burners. With freshly built ordinary recipient grids its medians were 179 versus 212 ms and 94 versus 198 ms. With the strongest retained ordinary index it instead took 191 versus 184 ms and 115 versus 102 ms. Cold component construction took a further 548--718 ms. Both woodpile stroke orders under this full-service policy had already burned every site before the second click. A direct empty-eligibility check answers both immediately. Separate full-state tests with an ordinary sparse scan of all positive stored heat also yielded mixed latency, even though that scan visited fewer sites.

We inspected the released grid solver of \citet{shen2026} as an independent comparison. Its wood temperature advances using occupied face neighbors and local heat loss. Its occupancy update reads the solid temperature and a fixed noise field. The gas reads the wood response. For a disposable wood-ignition query this one-way dependency lets an ordinary solver omit gas altogether. On an archived $128^3$ Firewoods run only 5,542 cells were occupied, 5,322 in one connected component. These facts limit the scope of the two-way material deferral claim.
